% GENERATED FILE. Edit canonical lessons, then run npm run generate:books.
% canonical-source-hash: fnv1a64:d35bf3cafa4382cd
% canonical-lessons: SA-C330-parampara, SA-C330-ritih, SA-C330-acarati, SA-C330-parva, SA-C330-tilakah

\chapter{Old Customs}
\label{ch:sa-old-customs}
\hlchaptermodality{330}

\begin{chapteropening}
\textbf{By the end of this chapter:} \emph{I can explain an old custom and how it is kept.}

Complete the last lesson of chapter 330: I can explain an old custom and how it is kept.
\end{chapteropening}

\section[\texorpdfstring{\sk{परम्परा} (paramparā)}{परम्परा (paramparā)}]{\sk{परम्परा} (paramparā) — a tradition}

\label{lesson:SA-C330-parampara}



\begin{quote}
Before the new one: say the Sanskrit for the waving of lamps, then the Sanskrit for prays for.
\end{quote}

\subsection*{You'll want to know: \sk{परम्परा}}
\textbf{\sk{परम्परा}} — \emph{paramparā} — \textquotedblleft{}a tradition\textquotedblright{}.

\begin{quote}
\sk{अस्माकं} \sk{ग्रामे} \sk{प्राचीना} \sk{परम्परा} \sk{अस्ति।}
\end{quote}

\begin{quote}
\emph{asmākaṃ grāme prācīnā paramparā asti.}
\end{quote}

\begin{quote}
Our village has an old tradition.
\end{quote}

A feminine noun.

\begin{cousinweb}
\textbf{\sk{परम्परा}} (\emph{paramparā}) is built from \emph{para}, \textquotedblleft{}next, other\textquotedblright{}, doubled with an \emph{m} between: one after another; MW gives \emph{paramparā} as an uninterrupted series, succession, and tradition (MW).
\end{cousinweb}

\begin{grammarlens}[title={What you've built}]
One more word for telling what happened.
\end{grammarlens}

\subsection*{Your turn}
\begin{itemize}
\raggedright
  \item \emph{Say it:} \emph{paramparā}
  \item \emph{Say it:} \emph{paramparā}, once more
  \item \emph{Say it:} \emph{asmākaṃ grāme prācīnā paramparā asti.} — Our village has an old tradition.
  \item \emph{Recall:} say the Sanskrit for the waving of lamps, then the Sanskrit for prays for, then say \emph{paramparā} again
\end{itemize}

\begin{tcolorbox}[breakable,colback=teal!4,colframe=teal!35!black,title={Before you move on}]
What does \sk{परम्परा} mean? (\textquotedblleft{}A tradition\textquotedblright{}.) Say it once more.
\end{tcolorbox}
\section[\texorpdfstring{\sk{रीतिः} (rītiḥ)}{रीतिः (rītiḥ)}]{\sk{रीतिः} (rītiḥ) — a custom; a way of doing things}

\label{lesson:SA-C330-ritih}



\begin{quote}
Before the new one: say the Sanskrit for prays for, then the Sanskrit for a tradition.
\end{quote}

\subsection*{You'll want to know: \sk{रीतिः}}
\textbf{\sk{रीतिः}} — \emph{rītiḥ} — \textquotedblleft{}a custom; a way of doing things\textquotedblright{}.

\begin{quote}
\sk{पितामहस्य} \sk{रीतिः} \sk{अद्य} \sk{अपि} \sk{जीवति।}
\end{quote}

\begin{quote}
\emph{pitāmahasya rītiḥ adya api jīvati.}
\end{quote}

\begin{quote}
Grandfather's custom still lives on today.
\end{quote}

A feminine noun.

\begin{cousinweb}
\textbf{\sk{रीतिः}} (\emph{rītiḥ}) MW's first senses are going, motion, a stream, and it also gives \emph{rīti} as a general way, usage, custom, method; MW gives no root in the entry, so its origin is left uncertain here (MW).
\end{cousinweb}

\begin{grammarlens}[title={What you've built}]
One more word for telling what happened.
\end{grammarlens}

\subsection*{Your turn}
\begin{itemize}
\raggedright
  \item \emph{Say it:} \emph{rītiḥ}
  \item \emph{Say it:} \emph{rītiḥ}, once more
  \item \emph{Say it:} \emph{pitāmahasya rītiḥ adya api jīvati.} — Grandfather's custom still lives on today.
  \item \emph{Recall:} say the Sanskrit for prays for, then the Sanskrit for a tradition, then say \emph{rītiḥ} again
\end{itemize}

\begin{tcolorbox}[breakable,colback=teal!4,colframe=teal!35!black,title={Before you move on}]
What does \sk{रीतिः} mean? (\textquotedblleft{}A custom; a way of doing things\textquotedblright{}.) Say it once more.
\end{tcolorbox}
\section[\texorpdfstring{\sk{आचरति} (ācarati)}{आचरति (ācarati)}]{\sk{आचरति} (ācarati) — observes; practises}

\label{lesson:SA-C330-acarati}



\begin{quote}
Before the new one: say the Sanskrit for a tradition, then the Sanskrit for a custom.
\end{quote}

\subsection*{You'll want to know: \sk{आचरति}}
\textbf{\sk{आचरति}} — \emph{ācarati} — \textquotedblleft{}observes; practises\textquotedblright{}.

\begin{quote}
\sk{वयम्} \sk{अद्य} \sk{अपि} \sk{प्राचीनां} \sk{रीतिम्} \sk{आचरामः।}
\end{quote}

\begin{quote}
\emph{vayam adya api prācīnāṃ rītim ācarāmaḥ.}
\end{quote}

\begin{quote}
We still observe the old custom today.
\end{quote}

Past (he …): \sk{सः} \sk{प्राचीनां} \sk{रीतिम्} \sk{आचरितवान्}, \emph{saḥ prācīnāṃ rītim ācaritavān}.

\begin{cousinweb}
\textbf{\sk{आचरति}} (\emph{ācarati}) is built from \emph{ā} and the root √car, \textquotedblleft{}to move\textquotedblright{}, as in \emph{carati}; MW gives \emph{ā-car} as to approach, to behave, and to act, practise, perform (MW).
\end{cousinweb}

\begin{grammarlens}[title={What you've built}]
One more word for telling what happened.
\end{grammarlens}

\subsection*{Your turn}
\begin{itemize}
\raggedright
  \item \emph{Say it:} \emph{ācarati}
  \item \emph{Say it:} \emph{ācarati}, once more
  \item \emph{Say it:} \emph{vayam adya api prācīnāṃ rītim ācarāmaḥ.} — We still observe the old custom today.
  \item \emph{Recall:} say the Sanskrit for a tradition, then the Sanskrit for a custom, then say \emph{ācarati} again
\end{itemize}

\begin{tcolorbox}[breakable,colback=teal!4,colframe=teal!35!black,title={Before you move on}]
What does \sk{आचरति} mean? (\textquotedblleft{}Observes; practises\textquotedblright{}.) Say it once more.
\end{tcolorbox}
\section[\texorpdfstring{\sk{पर्व} (parva)}{पर्व (parva)}]{\sk{पर्व} (parva) — a festival day}

\label{lesson:SA-C330-parva}



\begin{quote}
Before the new one: say the Sanskrit for a custom, then the Sanskrit for observes a custom.
\end{quote}

\subsection*{You'll want to know: \sk{पर्व}}
\textbf{\sk{पर्व}} — \emph{parva} — \textquotedblleft{}a festival day\textquotedblright{}.

\begin{quote}
\sk{अद्य} \sk{पर्व} \sk{अस्ति}, \sk{जनाः} \sk{नूतनानि} \sk{वस्त्राणि} \sk{धारयन्ति।}
\end{quote}

\begin{quote}
\emph{adya parva asti, janāḥ nūtanāni vastrāṇi dhārayanti.}
\end{quote}

\begin{quote}
Today is a festival day, and people wear new clothes.
\end{quote}

A neuter noun.

\begin{cousinweb}
\textbf{\sk{पर्व}} (\emph{parva}) is the nominative of \emph{parvan}, which MW gives first as a knot or joint, then a break, a section, and also a festival or holiday; MW gives no root, so its origin is left uncertain here (MW).
\end{cousinweb}

\begin{grammarlens}[title={What you've built}]
One more word for telling what happened.
\end{grammarlens}

\subsection*{Your turn}
\begin{itemize}
\raggedright
  \item \emph{Say it:} \emph{parva}
  \item \emph{Say it:} \emph{parva}, once more
  \item \emph{Say it:} \emph{adya parva asti, janāḥ nūtanāni vastrāṇi dhārayanti.} — Today is a festival day, and people wear new clothes.
  \item \emph{Recall:} say the Sanskrit for a custom, then the Sanskrit for observes a custom, then say \emph{parva} again
\end{itemize}

\begin{tcolorbox}[breakable,colback=teal!4,colframe=teal!35!black,title={Before you move on}]
What does \sk{पर्व} mean? (\textquotedblleft{}A festival day\textquotedblright{}.) Say it once more.
\end{tcolorbox}
\section[\texorpdfstring{\sk{तिलकः} (tilakaḥ)}{तिलकः (tilakaḥ)}]{\sk{तिलकः} (tilakaḥ) — a forehead mark; a tilaka}

\label{lesson:SA-C330-tilakah}



\begin{quote}
Before the new one: say the Sanskrit for observes a custom, then the Sanskrit for a festival day.
\end{quote}

\subsection*{You'll want to know: \sk{तिलकः}}
\textbf{\sk{तिलकः}} — \emph{tilakaḥ} — \textquotedblleft{}a forehead mark; a tilaka\textquotedblright{}.

\begin{quote}
\sk{माता} \sk{पुत्रस्य} \sk{ललाटे} \sk{तिलकं} \sk{करोति।}
\end{quote}

\begin{quote}
\emph{mātā putrasya lalāṭe tilakaṃ karoti.}
\end{quote}

\begin{quote}
Mother puts a mark on her son's forehead.
\end{quote}

A masculine noun.

\begin{cousinweb}
\textbf{\sk{तिलकः}} (\emph{tilakaḥ}) MW gives \emph{tilaka} as a mark on the forehead, of coloured earth, sandal or unguents, and as a freckle, without a derivation, so its origin is uncertain (MW).
\end{cousinweb}

\begin{grammarlens}[title={What you've built}]
One more word for telling what happened.
\end{grammarlens}

\subsection*{Your turn}
\begin{itemize}
\raggedright
  \item \emph{Say it:} \emph{tilakaḥ}
  \item \emph{Say it:} \emph{tilakaḥ}, once more
  \item \emph{Say it:} \emph{mātā putrasya lalāṭe tilakaṃ karoti.} — Mother puts a mark on her son's forehead.
  \item \emph{Recall:} say the Sanskrit for observes a custom, then the Sanskrit for a festival day, then say \emph{tilakaḥ} again
\end{itemize}

\begin{tcolorbox}[breakable,colback=teal!4,colframe=teal!35!black,title={Before you move on}]
What does \sk{तिलकः} mean? (\textquotedblleft{}A forehead mark; a tilaka\textquotedblright{}.) Say it once more.
\end{tcolorbox}
